\section{真题分类带练}

\subsection{集合}

\paragraph{1.}（2024 新课标Ⅰ卷）已知集合 $A = \{x \mid -5 < x^3 < 5\}$，$B = \{-3,-1,0,2,3\}$，则 $A \cap B =$
A. $\{-1,0\}$ \quad B. $\{2,3\}$ \quad C. $\{-3,-1,0\}$ \quad D. $\{-1,0,2\}$

\textbf{答案：A}

\subsection{复数}

\paragraph{2.}（2023 新课标Ⅰ卷）已知 $z = \dfrac{1-i}{2+2i}$，则 $z - \bar{z} =$
A. $-i$ \quad B. $i$ \quad C. $0$ \quad D. $1$

\textbf{答案：A}

\subsection{函数与导数}

\paragraph{3.}（2024 新课标Ⅰ卷）已知函数 $f(x) = \begin{cases} e^x + a, & x \le 0 \\ \ln(x+1) + x, & x > 0 \end{cases}$ 在 $\mathbb{R}$ 上单调递增，则 $a$ 的取值范围是
A. $(-\infty,0]$ \quad B. $[-1,0]$ \quad C. $[-1,+\infty)$ \quad D. $[0,+\infty)$

\textbf{答案：B}

\paragraph{4.}（2022 新课标Ⅰ卷）设 $a = 0.1e^{0.1}$，$b = \dfrac{1}{9}$，$c = -\ln 0.9$，则
A. $a < b < c$ \quad B. $c < b < a$ \quad C. $c < a < b$ \quad D. $a < c < b$

\textbf{答案：C}

\subsection{三角函数}

\paragraph{5.}（2024 新课标Ⅰ卷）已知 $\cos^2\alpha + 2\sin2\alpha = \dfrac15$，则 $\tan\alpha =$
A. $5\pm2\sqrt6$ \quad B. $10\pm2\sqrt{26}$ \quad C. $5\pm2\sqrt{26}$ \quad D. $10\pm2\sqrt6$

\textbf{答案：B}

\paragraph{6.}（2023 新课标Ⅰ卷）在 $\triangle ABC$ 中，$A+B=3C$，$2\sin(A-C)=\sin B$。
（1）求 $\sin A$；（2）设 $AB=5$，求 $AB$ 边上的高。

\textbf{答案}：（1）$\dfrac{3\sqrt{10}}{10}$；（2）$6$

\subsection{数列}

\paragraph{7.}（2022 新课标Ⅰ卷）已知数列 $\{a_n\}$ 满足 $a_1=1$，$a_{n+1} = \begin{cases}a_n+1, & n\text{ 为奇数} \\ a_n+3, & n\text{ 为偶数}\end{cases}$。
（1）记 $b_n = a_{2n}$，求 $\{b_n\}$ 的通项公式；
（2）求 $\{a_n\}$ 的前 20 项和。

\textbf{答案}：（1）$b_n = 4n-2$；（2）$S_{20}=300$

\subsection{立体几何}

\paragraph{8.}（2024 新课标Ⅰ卷）四棱锥 $P-ABCD$ 中，$PA\perp$ 底面 $ABCD$，$AD\perp CD$，$AD\parallel BC$，$PA=AD=CD=2$，$BC=3$。$E$ 为 $PD$ 中点。求平面 $PAB$ 与平面 $PCD$ 夹角的余弦值。

\textbf{答案}：$\dfrac{\sqrt{2}}{2}$

\subsection{概率统计}

\paragraph{9.}（2023 新课标Ⅰ卷）甲、乙两人投篮，每次由其中一人投篮。若命中则此人继续投篮，若未命中则换为对方投篮。甲命中概率 $0.6$，乙命中概率 $0.8$。第 1 次投篮的人是甲、乙的概率各为 $0.5$。求第 2 次投篮的人是乙的概率。

\textbf{答案}：$0.6$

\subsection{解析几何}

\paragraph{10.}（2024 新课标Ⅰ卷）椭圆 $E:\dfrac{x^2}{a^2}+\dfrac{y^2}{b^2}=1\;(a>b>0)$ 右焦点 $F$，点 $M(1,\frac32)$ 在 $E$ 上，$MF\perp x$ 轴。过 $F$ 的直线交 $E$ 于 $A,B$，$N$ 为 $AB$ 中点。求直线 $ON$ 斜率的最大值。

\textbf{答案}：（1）$\dfrac{x^2}{4}+\dfrac{y^2}{3}=1$；（2）$\dfrac{\sqrt{3}}{4}$

\subsection{多选综合}

\paragraph{11.}（2024 新课标Ⅰ卷）已知函数 $f(x)=x^3-3x^2+ax+b$，则
A. $f(x)$ 一定有两个极值点
B. 若有极大值点 $x_1$ 和极小值点 $x_2$，则 $x_1+x_2=2$
C. 若 $f(x)$ 在 $x=1$ 处取极值，则 $a=3$
D. 若 $f(x)$ 有三个零点，则 $a<3$

\textbf{答案：BC}

\subsection*{练习建议}

\[
\begin{array}{c|c|c}
\text{目标分数} & \text{重点突破} & \text{建议用时} \\ \hline
90\text{ 分} & \text{单选1-6、多选9-10、填空13-14、解答17-18} & \text{每天1h} \\
120\text{ 分} & \text{全部基础题 + 解答19-20、选填压轴部分} & \text{每天2h} \\
140+\text{ 分} & \text{全部题目 + 解答21-22} & \text{每天3h}
\end{array}
\]
